چون مجموعهٔ $S$ شمارا است، می‌توانیم تمام اعضای آن را به صورت یک دنبالهٔ مرتب از نقاط نام‌گذاری کنیم:
\[ S = \{p_1, p_2, p_3, \dots\} \]

مجموعه‌های $A$ و $B$ را در ابتدا تهی فرض می‌کنیم ($A = \emptyset$ و $B = \emptyset$). روند عضویت نقاط را گام‌به‌گام و براساس استقرار روی اندیس نقاط جلو می‌بریم.

برای هر گام $n = 1, 2, 3, \dots$، نقطهٔ $p_n = (x_n, y_n)$ را در نظر می‌گیریم:
\begin{itemize}
    \item فرض کنید $L_H^{(n)}$ خط افقی گذرنده از $p_n$ (یعنی خط $y = y_n$) باشد.
    \item فرض کنید $L_V^{(n)}$ خط عمودی گذرنده از $p_n$ (یعنی خط $x = x_n$) باشد.
\end{itemize}

در لحظه‌ای که می‌خواهیم تکلیف نقطهٔ $p_n$ را برای تخصیص به یکی از دو مجموعه مشخص کنیم، تعداد نقاط تخصیص‌یافته روی این خطوط تا قبل از این گام (یعنی از میان نقاط $p_1$ تا $p_{n-1}$) را به صورت زیر می‌شماریم:

$N_A(L_H^{(n)})$ تعداد نقاطی روی خط افقی $L_H^{(n)}$ است که به عضویت مجموعهٔ $A$ درآمده‌اند.

$N_B(L_V^{(n)})$ تعداد نقاطی روی خط عمودی $L_V^{(n)}$ است که به عضویت مجموعهٔ $B$ درآمده‌اند.



حالا نقطهٔ $p_n$ را بر اساس قانون تعادلی زیر به یکی از مجموعه‌ها می‌فرستیم:
\[
p_n \in 
\begin{cases} 
A & \text{if } \ \ \ \  N_A(L_H^{(n)}) \le N_B(L_V^{(n)}) \\ 
B & \text{if } \ \ \ \ =N_A(L_H^{(n)}) > N_B(L_V^{(n)}) 
\end{cases}
\]

\subsubsection*{بررسی صحت الگوریتم:}

 فرض کنید خط افقی $L_H: y = y_0$ وجود دارد که شامل {بی‌نهایت} نقطه از مجموعهٔ $A$ است. این نقاط را به ترتیب ظهورشان در دنبالهٔ اصلی $S$ با $a_1, a_2, a_3, \dots$ نشان می‌دهیم.

وقتی نقطهٔ $a_k$ بررسی می‌شود، چون به مجموعهٔ $A$ رفته است، طبق الگوریتم حتماً باید شرط زیر برقرار بوده باشد:
\[ N_A(L_H) \le N_B(L_V^{(k)}) \]
که در آن $L_V^{(k)}$ خط عمودی گذرنده از $a_k$ است.

از طرفی می‌دانیم $a_k$ امین نقطه‌ای روی این خط افقی است که به مجموعهٔ $A$ می‌رود؛ بنابراین در لحظهٔ بررسی آن، تعداد نقاط قبلی در $A$ روی این خط دقیقاً برابر $N_A(L_H) = k - 1$ بوده است. با جایگذاری این مقدار در نامساوی بالا خواهیم داشت:
\[ N_B(L_V^{(k)}) \ge k - 1 \]

این رابطه به این معنی است که به ازای هر $k \ge 1$، خط عمودی گذرنده از نقطهٔ $a_k$ باید قبل از بررسی $a_k$، حداقل شامل $k-1$ نقطهٔ متمایز متعلق به مجموعهٔ $B$ باشد.

وقتی $k \to \infty$ میل می‌کند، تعداد نقاط مورد نیاز در مجموعهٔ $B$ روی خطوط عمودی به صورت نامحدود رشد می‌کند. اما کل نقاط پردازش‌شده تا گام بررسی $a_k$ متناهی است (برابر با اندیس $a_k$ در مجموعهٔ $S$). این تناقض نشان می‌دهد که فرض اولیه‌مان غلط بوده و هیچ خط افقی نمی‌تواند شامل بی‌نهایت نقطه از $A$ باشد.

